%%%%%%%%%%%%%%%%%%%%%%%%%%%%%%%%%%%%%%%%%%%%%%%%%%%%%%%%%%%%%%%%%%%%%%%%%%%%%%%%
% referencial.tex
%
% Capítulo de referencial teórico da monografia.
%%%%%%%%%%%%%%%%%%%%%%%%%%%%%%%%%%%%%%%%%%%%%%%%%%%%%%%%%%%%%%%%%%%%%%%%%%%%%%%%

\chapter{Revisão Bibliográfica}
\label{sec:revbib}

Este capítulo tem como objetivo apresentar os fundamentos teóricos que englobam a proposta de
automação do fluxo pedagógico no sistema Educa EMES. A revisão parte de
conceitos de Sistemas de Informação, gestão educacional e automação de
processos, avançando para fundamentos de engenharia de software, arquitetura de
aplicações, padrões de design e elementos tecnológicos específicos, como
Laravel, MySQL, Livewire, Blade, Filament, comandos e processamento assíncrono.
A construção desse percurso teórico busca sustentar, em nível conceitual e
técnico, as decisões arquiteturais, de modelagem e de implementação discutidas
nos capítulos seguintes.

\section{Sistemas de Informação}
\label{sec:revbib:sistemas-informacao}

Sistemas de Informação podem ser compreendidos como arranjos que articulam
pessoas, processos organizacionais, dados e recursos tecnológicos para apoiar a
operação, o controle e a tomada de decisão nas organizações. Laudon e Laudon
\cite{LaudonLaudon2022} destacam que esses sistemas não se limitam à dimensão
técnica, pois dependem da integração entre tecnologia, processos de negócio e
necessidades organizacionais.

No contexto deste trabalho, essa perspectiva é relevante porque o Educa EMES
não é analisado apenas como aplicação computacional, mas como parte de um fluxo
institucional que envolve planejamento pedagógico, inscrições, frequência,
certificação, comunicação e acompanhamento de resultados. Assim, a automação
proposta precisa considerar tanto a arquitetura técnica quanto a aderência aos
processos reais executados pela EMES.

\section{Sistemas de Informação para Gestão Educacional}
\label{sec:revbib:gestao-educacional}

Sistemas de informação voltados à gestão educacional apoiam a coleta,
integração, processamento e disseminação de dados necessários ao planejamento,
monitoramento e administração de atividades educacionais. A UNESCO
\cite{UNESCOEMIS2026} define sistemas de informação de gestão educacional como
estruturas voltadas a apoiar a tomada de decisão e a gestão em diferentes
níveis do sistema educacional.

Embora o Educa EMES esteja inserido em uma escola da magistratura, e não em uma
rede educacional ampla, sua função institucional aproxima-se dessa lógica ao
centralizar dados e rotinas relacionados a cursos, participantes, instrutores,
frequências, certificados, documentos e relatórios. Essa aproximação reforça a
necessidade de padronização, rastreabilidade e consistência das informações ao
longo do fluxo pedagógico.

\section{Arquitetura de Software e Engenharia de Sistemas}
\label{sec:revbib:engsoft}

A engenharia de software oferece fundamentos importantes para compreender que a
qualidade de um sistema não depende apenas de sua existência ou de sua
capacidade de executar funções isoladas, mas da forma como seus elementos são
modelados, estruturados e organizados ao longo do desenvolvimento.
Sommerville \cite{Sommerville2011} destaca que sistemas de software devem ser
construídos com atenção a atributos como manutenibilidade, confiabilidade,
eficiência e adequação ao contexto de uso. De modo semelhante, Pressman e
Maxim \cite{Pressman2014} observam que decisões de arquitetura, modelagem e
projeto influenciam diretamente a capacidade de evolução da solução e sua
aderência aos requisitos do domínio.

A separação clara de responsabilidades, redução de acoplamento e organização
modular de componentes constituem princípios essenciais para sistemas capazes
de evoluir sem comprometer a estabilidade. Em contextos onde a automação de
processos é o objetivo central, uma arquitetura bem definida torna-se ainda
mais relevante, pois precisa sustentar não apenas a execução de funcionalidades,
mas também a extensibilidade da solução sem modificação de código existente
\cite{Sommerville2011,Pressman2014}.

Para sistemas educacionais integrados, essa discussão é particularmente crítica.
Um sistema deve representar com coerência elementos como cursos, inscrições,
critérios de aprovação, registros de frequência, certificação e comunicação com
participantes. Quando essa modelagem é insuficiente ou pouco aderente ao
processo real, a automação tende a ser parcial, frágil ou inconsistente.


\section{PHP}
\label{sec:revbib:php}
PHP (Hypertext Preprocessor) é uma linguagem de script amplamente utilizada para desenvolvimento web.
Criada por Rasmus Lerdorf em 1994, PHP evoluiu para se tornar uma das linguagens mais populares para
construção de aplicações web dinâmicas \cite{PHP2024}. PHP é uma linguagem de propósito geral, interpretada,
com tipagem dinâmica e suporte a múltiplos paradigmas de programação (procedural, orientada a objetos, funcional).
Sua sintaxe é inspirada em C, Java e Perl, o que facilita a adoção por desenvolvedores familiarizados com essas linguagens.


\section{Laravel e Model-View-Controller}
\label{sec:revbib:mvc}

Os frameworks web modernos implementam o padrão de arquitetura
\textit{Model-View-Controller} (MVC), que divide a aplicação em três camadas:
 Model (dados e lógica de negócio), View (visualização e apresentação) e Controller
(orquestração de requisições e respostas). Essa separação é fundamental para a organização do codigo,
permitindo que cada camada evolua de maneira independente, facilitando manutenção,
testabilidade e reutilização de código \cite{Sommerville2011}.

Laravel \cite{Laravel2024} é um framework PHP que utiliza do MVC em sua arquitetura
através de abstrações de alto nível. Stauffer \cite{CodingStyle2019} descreve
Laravel como uma plataforma que combina elegância sintática com potência
funcional, fornecendo ferramentas integradas para roteamento, persistência de
dados (via Eloquent ORM), validação, autenticação e autorização. A filosofia
do framework prioriza a experiência do desenvolvedor, se apresentando como uma solução rapidamente escalável, modular e fácil de manter,
permitindo focagem nas regras de negócio.

Um aspecto central do Laravel é o \textit{Eloquent Object-Relational Mapping}
(ORM) \cite{LaravelDocs2024}, responsável por abstrair o acesso ao banco de dados relacional,
permitindo que desenvolvedores trabalhem diretamente com objetos PHP, sem a necessidade de escrever SQL diretamente.
Isso facilita modelagem, construção de relacionamentos e consultas. Além disso, Eloquent suporta eventos
de ciclo de vida (creating, created, updating, deleting, etc.), permitindo logicas customizadas durante o ciclo de vida dos modelos.

\section{Banco de Dados Relacional e MySQL}
\label{sec:revbib:bd}

O modelo relacional de dados, formalizado por Date e Codd, permanece como
paradigma com maior adesão para armazenamento estruturado.Os bancos relacionais organizam as informações em tabelas
que podem portar relações com linhas e colunas, permitindo consultas via Structured Query Language (SQL).

MySQL \cite{MySQL2024} é um sistema de gerenciamento de banco de dados (SGBD)
relacional open source amplamente utilizado e difundido. Oferece
suporte a transactions, constraints para integridade referencial, triggers, procedimentos, além de índices para otimização de consultas.
Para sistemas de gestão educacional que dependem de confiabilidade, consistênciae conformidade com normas de integridade,
MySQL é uma boa escolha.

A modelagem adequada do esquema de dados é fundamental. Boas práticas incluem
normalização com suas formas normais, definição clara de relacionamentos,
uso apropriado de tipos de dados e criação de índices em colunas de busca
frequente. Quando utilizadas de maneira oportuna e correta, podemos garantir
que o banco
de dados além de permitir consultas eficientes será integro.

\section{Padrões de Design e Arquitetura}
\label{sec:revbib:patterns}

Padrões de design são soluções reutilizáveis para problemas recorrentes em
design de software. Gamma e colaboradores \cite{GammaEtAl1994}, através do
seminal ``Gang of Four'', categorizam padrões em três grupos: criacionais
(construção de objetos), estruturais (composição de objetos) e comportamentais
(comunicação entre objetos).

\subsection{Factory Pattern}
O padrão Factory encapsula a criação de objetos, permitindo que o cliente não
conheça as subclasses concretas. Isso é particularmente útil quando a
instanciação depende de decisões em tempo de execução ou quando se deseja
extensibilidade sem modificação de código cliente. Em contextos de automação,
factory facilita a criação dinâmica de estratégias baseada em configuração.

\subsection{Strategy Pattern}
O Strategy define uma família de algoritmos, encapsula cada um e os torna
intercambiáveis. Permite que o algoritmo varie independentemente do cliente
que o utiliza. Em sistemas de automação de fluxos, cada etapa pode possui
múltiplas estratégias de execução (elegibilidade, validação, ação) que são
selecionadas em tempo de execução conforme configuração ou estado do processo.

\subsection{Builder Pattern}
O Builder separa a construção de um objeto complexo de sua representação,
permitindo que o mesmo processo construa representações diferentes. Útil para
estruturas que possuem múltiplos componentes opcionais ou que precisam ser
construídas gradualmente. Em automação, builder é apropriado para construção
de estruturas de análise, resultado e diagnóstico com componentes variáveis.

\subsection{Service Layer}
O padrão Service Layer encapsula lógica de negócio em serviços que atuam como
fachada entre controllers (orquestração de requisição) e repositórios (acesso
a dados). Martin Fowler \cite{ServiceLayerPattern2024} identifica que esse
padrão favorece testabilidade, reutilização de lógica entre múltiplos
endpoints e clareza no fluxo de execução. Serviços podem ser compostos,
permitindo que lógica complexa seja construída a partir de componentes menores.

\section{Templating e Rasterização de Views}
\label{sec:revbib:templates}

Em arquiteturas web MVC, a camada de apresentação precisa ser renderizada
dinamicamente na base de dados e estado da aplicação. Blade
\cite{Blade2024} é o mecanismo de templating do Laravel que permite adicionar
lógica controlada (condicionais, loops, includes) em templates HTML sem perder
legibilidade. Blade compila templates para PHP em cache, fazendo com que não
haja overhead de parsing em requisições subsequentes.

O uso de templates bem estruturados favorece separação entre lógica de
apresentação e lógica de negócio, reduz duplicação de HTML e facilita
manutenção de interfaces quando regras de apresentação mudam. Em interfaces
administrativas complexas, templates reutilizáveis tornam-se essenciais para
eficiência de desenvolvimento.

\section{Componentes Reativos e Livewire}
\label{sec:revbib:components}

Aplicações web modernas frequentemente exigem interatividade sem recarregar a
página, historicamente alcançado através de JavaScript (cliente) e requisições
AJAX (assíncrono). Livewire \cite{Livewire2024} é uma biblioteca PHP que traz
reatividade ``full-stack'', permitindo que componentes PHP reajam a eventos do
usuário sem necessidade explícita de JavaScript.

Livewire possibilita que desenvolvedores escrevam componentes com estado
(properties) que, ao serem modificados, automaticamente atualizam a view sem
recarregar a página. Isso reduz a distância entre servidor e cliente,
eliminando a necessidade de APIs REST intermediárias para operações simples,
ao mesmo tempo que mantém a lógica no servidor onde pode ser testada e
auditada.

Para interfaces de automação onde operações são disparadas por ações de usuário
(iniciar um passo, confirmar eligibilidade, capturar diagnósticos), Livewire
oferece abstração apropriada reduzindo complexidade de JavaScript enquanto
mantém responsividade.

\section{Painéis Administrativos: Filament}
\label{sec:revbib:admin}

Sistemas integrados de gestão frequentemente requerem interfaces administrativas
robustas para configuração, monitoramento e operação. Filament
\cite{Filament2024,FilamentAdminPanel2024} é uma coleção de componentes
full-stack para Laravel que acelera a construção de painéis administrativos
profissionais.

Filament fornece abstrações de alto nível para CRUD (Create, Read, Update,
Delete), tabelas interativas, filtros, ações em lote e customização visual.
Reduz significativamente tempo de desenvolvimento de interfaces administrativas
mantendo qualidade e consistência visual. Para sistemas educacionais que
necessitam de várias interfaces específicas de diferentes usuários
(administradores de sistema, gestores educacionais, coordenadores, instrutores),
Filament possibilita construção rápida sem perder controle sobre modelagem e
lógica de negócio \cite{FilamentAdminPanel2024}.

\section{Automação de Processos e Orquestração}
\label{sec:revbib:automação}

A automação de processos organizacionais refere-se ao uso de mecanismos
sistêmicos para executar, encadear ou apoiar atividades de maneira padronizada,
reduzindo a dependência de intervenções manuais e aumentando a previsibilidade
da operação. Em sistemas de informação, a automação não corresponde apenas à
substituição de ações humanas por código, mas à tradução de regras, estados,
eventos e critérios de decisão para mecanismos computacionais coerentes com o
processo real.

Entre benefícios estão redução de falhas operacionais, rastreabilidade ampliada,
padronização de rotinas e ganhos de eficiência. Porém, processos mal modelados
ou pouco aderentes ao contexto podem produzir novos gargalos. A automação do
fluxo pedagógico, portanto, não é meramente implementação de código, mas
atividade que envolve engenharia de processos e tradução adequada das regras
institucionais para a estrutura do sistema.

\section{Gestão de Processos de Negócio e BPMN}
\label{sec:revbib:bpm-bpmn}

A discussão sobre automação aproxima-se do campo de \textit{Business Process
Management} (BPM), no qual processos são tratados como estruturas
organizáveis, modeláveis, executáveis e passíveis de melhoria contínua
\cite{Weske2024,DumasEtAl2018}. Nessa perspectiva, a análise do processo não se
restringe a funções isoladas, mas observa o encadeamento das atividades, as
regras de negócio, os responsáveis envolvidos e as dependências entre eventos.

No contexto deste trabalho, a abordagem de BPM contribui para compreender o
fluxo pedagógico da EMES como uma sequência de etapas interdependentes, desde o
planejamento e a inscrição até a frequência, aprovação, certificação e geração
de relatórios. Essa compreensão é necessária para que a automação não apenas
execute tarefas pontuais, mas preserve coerência entre estados, critérios de
transição, pendências e possibilidades de intervenção manual.

BPMN é uma notação padronizada para representação de processos de negócio,
permitindo documentar fluxos atuais, projetar fluxos futuros e comunicar etapas,
eventos e decisões de maneira compreensível entre perfis técnicos e
organizacionais \cite{OMG2014}. Por esse motivo, a modelagem de processos
assume papel importante na análise do fluxo pedagógico, pois ajuda a explicitar
dependências, gargalos e pontos em que a automação pode ser aplicada.

\section{Processamento Assíncrono e Job Queues}
\label{sec:revbib:async}

Em sistemas onde operações podem ser computacionalmente intensivas ou precisam
executar em background sem bloquear a requisição HTTP, processamento assíncrono
torna-se essencial. Laravel oferece mecanismos integrados para isso: Queues
\cite{LaravelQueues2024} e Tasks Scheduling \cite{LaravelScheduler2024}.

Queues permitem que tarefas sejam adicionadas a fila (em memória ou persistidas
em redis/banco de dados) e processadas por workers em background. Isso
desacopla o tempo de enfileiramento do tempo de processamento, permitindo que
requisições HTTP retornem rapidamente ao usuário enquanto o trabalho pesado
ocorre posteriormente. Scheduling, por sua vez, permite que comandos sejam
executados periodicamente (ex: a cada minuto, hora ou dia) sem necessidade de
cron jobs scripts externos.

Para automação de fluxos pedagógicos, esses mecanismos são relevantes quando
passos de automação demandam processamento pesado, notificações múltiplas ou
execução em horários específicos. Por exemplo, processamento de certificações
em lote, envio de notificações, geração de relatórios ou sincronização com
sistemas externos pode ser delegado a filas, deixando endpoints HTTP livres
para requisições interativas \cite{LaravelQueues2024}.

\section{Commands e Artisan Console}
\label{sec:revbib:commands}

Artisan \cite{LaravelCommands2024} é a interface de linha de comando do Laravel,
fornecendo uma série de comandos built-in para gerenciamento de aplicação
(criar controllers, modelos, migrations, etc.) e um framework para criação de
comandos customizados.

Commands customizados encapsulam rotinas operacionais que precisam ser
executadas fora do contexto de uma requisição HTTP. Exempos incluem
importação/exportação de dados, sincronização com sistemas externos, limpeza
de cache, processamento em lote e orquestração de fluxos de automação.

Quando um Command é despachado (seja manualmente via CLI ou via Scheduler),
ele executa em um contexto isolado com acesso total ao container de serviços,
banco de dados e demais infraestrutura da aplicação. Para automação de fluxos,
Commands atuam como ponto de entrada principal: buscar entidades elegíveis
para processamento, aplicar regras de elegibilidade, disparar ações e
coordenar orquestração de passos.

\section{Arquitetura Orientada a Eventos}
\label{sec:revbib:events}

Em sistemas complexos onde múltiplos componentes precisam reagir ao mesmo evento
sem acoplamento direto, o padrão Observer (ou Event-Driven) oferece solução
limpa. Martin Fowler \cite{EventDrivenArchitecture2024} descreve Event Sourcing
como abordagem em que eventos significativos do domínio são capturados,
persistidos e podem ser replicados para múltiplos subscribers.

Em Laravel, eventos podem ser disparados quando ocorrem transições significativas
no domínio (ex: inscrição criada, aprovação concedida, certificação gerada).
Listeners podem estar registrados para disparar ações adicionais (enviar email,
atualizar dashboard, disparar outro processo) sem que o código que dispara o
evento conheca os listeners \cite{LaravelDocs2024}.

Para automação de fluxos pedagógicos, essa abordagem é particularmente apropriada:
quando uma entidade transiciona de estado (ex: curso muda de ``pendente
aprovação'' para ``aprovado''), eventos podem ser disparados para notificar
interessados, iniciar novos processos ou atualizar registros relacionados, sem
criar acoplamento excessivo entre diferentes partes do sistema.

\section{Trabalhos Relacionados}
\label{sec:trabalhos-relacionados}

Esta seção apresenta trabalhos relacionados à modelagem, melhoria e automação de processos
acadêmicos e administrativos em instituições de ensino ou em contextos organizacionais 
próximos. A análise desses estudos tem como objetivo posicionar a proposta deste trabalho 
em relação a iniciativas que utilizam Business Process Management (BPM), Business Process 
Model and Notation (BPMN), levantamento de requisitos, prototipação e automação de fluxos 
institucionais. A comparação concentra-se especialmente nas etapas de análise do processo 
atual, identificação de gargalos, modelagem de processos e definição de uma solução de software.

\subsection{Modelagem de processos acadêmicos com BPMN}
\label{sec:trabalhos-relacionados:bpmn-academico}

Fontes, Santos e Libório \cite{FontesSantosLiborio2020} aplicaram BPMN na melhoria de processos 
acadêmicos do Instituto Federal de Sergipe (IFS), com foco em processos diretamente percebidos 
pelos discentes, como matrícula e trancamento de matrícula. O estudo identificou problemas de 
clareza, burocracia e dependência de procedimentos manuais, propondo modelos futuros orientados 
à automação, como notificação individualizada por e-mail, gestão automática de lista de espera e 
início de solicitações diretamente pelo sistema acadêmico. A relevância desse trabalho para esta 
pesquisa está na demonstração de que processos acadêmicos recorrentes podem ser analisados e 
redesenhados com apoio de BPMN, resultando em propostas capazes de reduzir tarefas manuais e 
melhorar a experiência dos usuários.

Martins e Pinho \cite{MartinsPinho2021} investigaram melhorias no setor de Registro Acadêmico 
de uma universidade federal utilizando a técnica BPMN. O estudo analisou processos como suspensão 
de curso, desligamento e colação de grau extraordinária, evidenciando falhas estruturais, 
dependência de formulários físicos e baixa utilização do sistema institucional existente. 
A principal contribuição do trabalho está na defesa de que o diagnóstico do processo atual 
deve anteceder qualquer tentativa de automação, evitando que fluxos ineficientes sejam simplesmente 
digitalizados. Essa perspectiva é especialmente relevante para o presente trabalho, pois a automação 
do fluxo pedagógico do Educa EMES também parte da análise de gargalos e da estrutura existente antes 
da proposição de mecanismos automatizados.

Braga e Aires \cite{BragaAires2024} realizaram a modelagem de processos em uma unidade de uma instituição de ensino superior, demonstrando o uso de BPM como instrumento de compreensão e melhoria de processos internos. Embora o estudo não tenha como foco principal a construção de uma solução computacional, ele reforça a importância da modelagem para tornar processos organizacionais explícitos, identificar falhas de comunicação e propor ajustes gerenciais. Em comparação com este trabalho, sua contribuição está mais associada à melhoria procedimental, enquanto a proposta do Educa EMES busca avançar para uma modelagem técnica integrada a um sistema educacional já existente.

\subsection{Da modelagem de processos ao projeto de software}
\label{sec:trabalhos-relacionados:modelagem-software}

Lima \cite{Lima2026} desenvolveu um trabalho voltado à otimização do processo de solicitação e matrícula em Trabalho de Conclusão de Curso (TCC) na Universidade Federal do Ceará, integrando BPMN à Engenharia de Software. A pesquisa mapeou o processo atual (AS-IS), identificou gargalos, elaborou um processo futuro (TO-BE), especificou requisitos em formato de histórias de usuário e desenvolveu um protótipo de alta fidelidade. Esse estudo é um dos mais próximos metodologicamente deste trabalho, pois utiliza a modelagem de processos como base para derivar requisitos e projetar uma solução de software. A diferença principal está no escopo e no estágio da solução: enquanto o trabalho de Lima concentra-se em um fluxo acadêmico específico e em um protótipo, o presente trabalho analisa um fluxo pedagógico mais amplo, envolvendo inscrição, frequência, aprovação, certificação e relatórios dentro de um sistema institucional já em produção.

Saraiva Júnior et al. \cite{SaraivaJuniorEtAl2025} apresentaram uma experiência de integração entre modelagem de processos de negócio e Engenharia de Requisitos no contexto da Fundação Cearense de Apoio ao Desenvolvimento Científico e Tecnológico (Funcap). O estudo utilizou modelos AS-IS e TO-BE para organizar processos institucionais e derivar requisitos de software, documentados por meio de histórias de usuário e mecanismos de rastreabilidade. Apesar de não estar restrito ao domínio educacional, o trabalho contribui para esta pesquisa ao evidenciar a utilidade da modelagem de processos como artefato intermediário entre a compreensão do domínio organizacional e a especificação de soluções computacionais.

\subsection{Automação de processos e workflows em instituições públicas e de ensino}
\label{sec:trabalhos-relacionados:automacao-workflows}

Galimberti, De Lucca e Ramos \cite{GalimbertiLuccaRamos2024} apresentaram uma experiência de gestão e automação de processos em uma universidade pública, discutindo a institucionalização de práticas de BPM e automação em um contexto governamental. O estudo é relevante por evidenciar que a automação de processos em instituições públicas exige mais do que a adoção pontual de ferramentas, demandando governança, priorização de processos, acompanhamento e organização metodológica. Essa perspectiva contribui para fundamentar a proposta deste trabalho, uma vez que a automação do fluxo pedagógico da EMES é tratada como evolução estruturada de um sistema real, e não como funcionalidade isolada.

Greavu-Șerban \cite{GreavuSerban2015} analisou a automação de processos administrativos em instituições de ensino superior na Romênia, com uso de workflows, BPMN, assinatura digital e plataformas corporativas. A pesquisa demonstra que processos acadêmico-administrativos, especialmente aqueles relacionados a documentos e certificações, tendem a apresentar alto grau de burocracia e podem se beneficiar da automação baseada em fluxos. A relação com este trabalho ocorre principalmente nas etapas de certificação, comunicação e disponibilização de documentos acadêmicos, que também compõem o fluxo pedagógico do Educa EMES.

Adedapo et al. \cite{AdedapoEtAl2023} propuseram um sistema de gerenciamento de workflow para um departamento universitário, com objetivo de centralizar comunicação, automatizar tarefas e melhorar a produtividade em um ambiente de ensino superior. Embora a solução utilize tecnologias distintas das empregadas no Educa EMES, o trabalho reforça a pertinência de sistemas orientados a workflow para lidar com a complexidade de processos acadêmicos e administrativos. Sua contribuição para esta pesquisa está na aproximação entre gestão acadêmica, automação de tarefas e acompanhamento de fluxos institucionais.

\subsection{Síntese comparativa}
\label{sec:trabalhos-relacionados:sintese}

A análise dos trabalhos relacionados demonstra que a modelagem de processos tem sido utilizada em instituições de ensino e organizações públicas principalmente para compreender fluxos existentes, identificar gargalos, propor melhorias e, em alguns casos, derivar requisitos de software. Parte dos estudos concentra-se na documentação e melhoria organizacional, enquanto outros avançam para a construção de protótipos ou sistemas de apoio. No contexto deste trabalho, a contribuição está na combinação entre análise de processo pedagógico, análise estrutural de um sistema existente e modelagem de uma solução de automação integrada ao Educa EMES.

\begin{table}[H]
\centering
\caption{Comparação entre trabalhos relacionados e este trabalho}
\label{tab:comparacao-trabalhos-relacionados}
\begin{tabular}{p{3.2cm} p{3.1cm} p{3.2cm} p{3.9cm}}
\hline
\textbf{Trabalho} & \textbf{Contexto} & \textbf{Foco principal} & \textbf{Relação com este trabalho} \\
\hline
Fontes, Santos e Libório \cite{FontesSantosLiborio2020} & Instituto Federal de Sergipe & Melhoria de processos acadêmicos com BPMN & Relaciona-se à automação de processos acadêmicos e redução de procedimentos manuais. \\
\hline
Martins e Pinho \cite{MartinsPinho2021} & Universidade federal & Registro acadêmico e diagnóstico de gargalos & Reforça a necessidade de analisar o processo atual antes de automatizar. \\
\hline
Braga e Aires \cite{BragaAires2024} & Instituição de ensino superior & Modelagem de processo interno & Contribui para justificar BPM como instrumento de compreensão e melhoria organizacional. \\
\hline
Lima \cite{Lima2026} & Universidade Federal do Ceará & Matrícula em TCC com BPMN, requisitos e protótipo & Aproxima-se metodologicamente por ligar AS-IS, TO-BE, requisitos e solução de software. \\
\hline
Saraiva Júnior et al. \cite{SaraivaJuniorEtAl2025} & Gestão pública & BPMN e Engenharia de Requisitos & Relaciona modelagem de processos à derivação de requisitos de software. \\
\hline
Galimberti, De Lucca e Ramos \cite{GalimbertiLuccaRamos2024} & Universidade pública & Gestão e automação de processos & Fundamenta a automação institucional em contexto público. \\
\hline
Greavu-Șerban \cite{GreavuSerban2015} & Ensino superior internacional & Automação administrativa com workflow e BPMN & Relaciona-se a processos acadêmico-administrativos, documentos e certificação. \\
\hline
Adedapo et al. \cite{AdedapoEtAl2023} & Departamento universitário & Sistema de workflow para ensino superior & Reforça o uso de sistemas orientados a workflow em processos acadêmicos. \\
\hline
Este trabalho & Escola da Magistratura do Espírito Santo & Proposta de automação no fluxo pedagógico da EMES & Diferencia-se por modelar uma evolução integrada a um sistema educacional real em produção. \\
\hline
\end{tabular}
\end{table}

De modo geral, os trabalhos analisados confirmam a relevância de BPM, BPMN, workflows e Engenharia de Requisitos para a melhoria de processos acadêmicos e administrativos. Entretanto, o presente trabalho diferencia-se por ter como objeto uma plataforma educacional institucional já existente, utilizada em ambiente real, e por propor a automação de um fluxo pedagógico mais abrangente, envolvendo diferentes etapas do ciclo de vida de cursos e eventos.
